\subsection{Related frameworks and decision-support applications}\label{sec:related}

Recent research increasingly treats decision support as a \emph{hybrid} process rather than as the output of a single isolated model. In practice, decision quality depends on how automated models, human expertise, uncertainty handling, and contextual interpretation are combined within an explicit workflow. For the present discussion, it is useful to distinguish between two neighboring literatures: (i) collaborative decision frameworks, which specify how humans and AI systems share authority under uncertainty, and (ii) application-oriented decision-support systems, which implement concrete fuzzy, MCDM, or hybrid ML pipelines for particular domains \citep{Tariq2024A2C,Deliu2026HAIC,Abdulgader2018Fuzzy,Sahraei2022Flood}.

At the framework level, the strongest direct analogue is the A2C framework of \citet{Tariq2024A2C}, which was explicitly designed to support automated, augmented, and collaborative decision modes within human--AI teams. A2C builds on rejection learning and learning-to-defer, enabling AI components to abstain in ambiguous cases and defer to human experts, while also supporting collaborative exploration when neither side can resolve uncertainty alone \citep{Tariq2024A2C}. Related work on hybrid intelligence likewise emphasizes that authority distribution and human interpretive capacity should be modeled separately: in project and innovation settings, the level of AI integration and the cognitive flexibility of human decision-makers jointly shape decision quality, accountability, and resilience \citep{Deliu2026HAIC}. In the strategic-management literature, the same pattern appears as a human--technology symbiosis in which AI contributes large-scale analytical and predictive capability, but humans retain interpretive and ethical authority over final choices \citep{Telaumbanua2025Symbiosis}.

A neighboring but distinct line of work studies cooperation among multiple AI agents, especially large-language-model agents. Recent surveys note that extending single-agent reinforcement learning to multi-agent settings is non-trivial because coordination, communication, and shared objectives must be learned explicitly \citep{Sun2024MARL}. More recent multi-agent RL approaches for LLM collaboration argue that prompt-level coordination alone is fragile, because fixed agents can propagate conflicting or incorrect information, whereas coordination-aware optimization can improve robustness and efficiency \citep{Liu2025CoMLRL}. Trust-weighted consensus mechanisms further suggest that collaborative reasoning can be strengthened when contributions from more reliable agents receive greater weight during group decision formation \citep{Shakya2026Trust}. These studies are relevant because they formalize cooperation among decision-producing units, but their primary goal is agent coordination or reasoning performance rather than a reusable decision-support workflow that keeps predictive outputs, context modulation, and final prioritization analytically distinct.

At the application level, a substantial literature has developed transparent decision-support pipelines based on fuzzy MCDM and hybrid fuzzy--ML integration. Hybrid fuzzy MCDM models are commonly used to combine interdependent criteria analysis, pairwise weighting, and alternative ranking in complex industrial decisions \citep{Abdulgader2018Fuzzy}. Fuzzy DEMATEL-based designs are particularly relevant because they separate cause--effect relationships among criteria, explicitly represent expert judgments, and manage the uncertainty inherent in linguistic assessments \citep{Sahraei2022Flood}. More recent hybrid systems also integrate machine learning prediction and explainable AI with fuzzy MCDM. For example, \citet{Ozkurt2025Energy} combine fuzzy TOPSIS and fuzzy ELECTRE with ML prediction and XAI tools, showing how ranking models and predictive models can coexist in a transparent decision-support architecture. Similarly, \citet{Mikhaylov2026Energy} integrate DEMATEL with a neuro-behavioral and picture-fuzzy rough framework to support regulatory oversight, while explicitly evaluating transparency, accountability, governance, and reliability criteria.

These contributions show the value of combining data-driven evidence with expert-defined or context-sensitive criteria, but they are usually tailored to a specific method stack or domain application. In most cases, one finds either (a) an advanced aggregation model, (b) a hybrid prediction-and-ranking pipeline, or (c) a collaborative decision framework; much less common is an operational software architecture that separates multiple ADM outputs from a dedicated context layer and then recomputes final prioritization under alternative decision policies \citep{Tariq2024A2C,Abdulgader2018Fuzzy,Ozkurt2025Energy}. This distinction is central to the positioning of CADEMAS-ML.

Open-source software is another relevant axis because reproducibility strongly affects the practical value of decision-support research. In this area, \textsc{pyFDM} provides a modular Python environment for fuzzy decision-making under uncertainty, including normalization, defuzzification, distance measures, and weighting techniques \citep{Wieckowski2022pyFDM}. \textsc{PyIFDM} extends this line to intuitionistic fuzzy MCDA, offering multiple methods, normalization procedures, similarity measures, and visualization support for complex uncertain decision models \citep{Wieckowski2023PyIFDM,Wieckowski2024PyIFDMUpdate}. Complementarily, the \texttt{IFMCDM} R package supports intuitionistic fuzzy ranking with alternative reference-point and distance formulations, increasing flexibility for comparative decision analysis \citep{Roszkowska2024IFMCDM}. Computational support has also been proposed for intuitionistic fuzzy dimensional analysis, reinforcing the broader trend toward executable and auditable MCDM tooling rather than purely conceptual models \citep{PerezDominguez2024IFDA}.

\begin{table}[t]
\caption{Representative published works related to CADEMAS-ML.\label{tab:related}}
\small
\begin{adjustwidth}{-\extralength}{0cm}
\begin{tabularx}{\fulllength}{>{\raggedright\arraybackslash}p{2.8cm} >{\raggedright\arraybackslash}X p{1.9cm} p{1.9cm} >{\raggedright\arraybackslash}X}
\toprule
\textbf{Work} & \textbf{Decision-support focus} & \textbf{Multiple ADM units} & \textbf{Explicit context layer} & \textbf{Relation to CADEMAS-ML} \\
\midrule
\citet{Tariq2024A2C} & Human--AI collaborative decision framework with automated, augmented, and collaborative modes & Partly & Partly & Strong framework-level precedent, but not a configurable fuzzy/ML prioritization workflow \\
\citet{Deliu2026HAIC} & Governance framework for human--AI collaboration and cognitive flexibility in project decisions & No & Yes & Highlights authority and interpretive context, but not multi-ADM integration \\
\citet{Abdulgader2018Fuzzy} & Hybrid fuzzy DEMATEL--AHP--TOPSIS decision model & No & Partly & Rich hybrid MCDM pipeline, but not organized around cooperating ADM outputs \\
\citet{Sahraei2022Flood} & Fuzzy DEMATEL/BWM/ANP/COPRAS pipeline for spatial risk prioritization & No & Partly & Explicit causal weighting under uncertainty, but application-specific \\
\citet{Ozkurt2025Energy} & Hybrid fuzzy MCDM + ML + XAI for transparent ranking and prediction & Partly & Partly & Combines ranking and prediction, but lacks a general cooperative ADM architecture \\
\citet{Mikhaylov2026Energy} & DEMATEL-based fuzzy hybrid DSS with transparency and governance criteria & No & Yes & Strong context and oversight perspective, but method-specific \\
\midrule
CADEMAS-ML (this work) & Configurable prioritization with ML-based ADM units and fuzzy context scenarios & Yes & Yes & Operationalizes separation of ADM outputs, context evaluation, and final prioritization \\
\bottomrule
\end{tabularx}
\end{adjustwidth}
\end{table}

Table~\ref{tab:related} highlights two consistent patterns in the literature. First, published systems already provide strong ingredients for transparent decision support, including deferral-aware collaboration, fuzzy causal analysis, hybrid ranking pipelines, explainable prediction, and reusable software libraries \citep{Tariq2024A2C,Sahraei2022Flood,Ozkurt2025Energy,Wieckowski2022pyFDM}. Second, these ingredients are usually developed in isolation. What remains comparatively underdeveloped is a reusable implementation in which several decision-producing models are treated as cooperating ADM components, their outputs are integrated explicitly, contextual assessment is modeled as a separate layer, and final prioritization can be recomputed under alternative policies.

CADEMAS-ML is positioned to address that implementation gap for predictive-model settings. Its contribution is not merely the inclusion of fuzzy or ML techniques, nor only software availability, but the explicit operationalization of a decomposed decision process: multiple ADM units, a separate context-sensitive evaluation stage, and an auditable final prioritization mechanism that can be varied without collapsing these roles into a single monolithic model.